% !TeX root = ../papers/nc-p3.tex
\section{反転なしBell直線の全流入と二平面分解を閉極限に持つ枝の被覆}
\label{sec:noflip-line-fibers}
\begin{theorem}\label{thm:noflip-line-fibers}
半単純直線の $s\ne0$, $s^2\ne1$, $L=1$ を閉軌道とする。
全流入窓は標準Koszul左右AS正則で、閉軌道一類と
$\mathbb P(1^4,2^3)$ の次の有限群による商で尽きる。各非零軌道の半安定境界は閉軌道だけである。
最小周期は1または2で、非零類の最小周期1の軌跡の既約成分次元は
\[
 \begin{array}{c|l}
 s^4\ne1 & 4,4,3,3,3,3,2,2\\
 s^2=-1 & 4,4,3,3,2 .
 \end{array}
\]
閉軌道の基点同型は $s\sim-s\sim1/s$ である。基点を揃えた後、重み付き射影座標に関する下記の
有限行列の軌道比較が全ての添字保持同型を決定する。
\end{theorem}
\begin{proof}[全正座標と同型作用]
$P_{abab}=1$ ($a=b$) または $s$ ($a\ne b$) の通常基底で六成分を再構成する。
各成分の四支持語に等しい正重みを与えれば一文字周辺が一様で、一位置だけ異なる二語もないので
この窓は多重安定である。指定した二平面の法線文字を0,1、他を2,3とし、その語中の法線数を $n$ とする。
\path{computations/verify_noflip_line_fibers.py} は $\mathbb Q[s,s^{-1}]$ 上で各切断の
左右核の積を取り、$n=3,4$ の接方程式の核次元が20,3で一定であることをSmith形で確認する。
使う核は単位座標を含み、分母は $s$ の冪だけである。背景の四純関係が単項可逆で混合ブロックも
可逆なので核の次元は特殊化で保たれる。せん断16列にも $s(s^2-1)$ の因子のみを持つ
非零な16次小行列式がある。残る重み1の四方向の各平方・混合積の両Schur右辺を
四法線係数の行列に加えても、一般体で階数は13のままである。
定階数の束の余核への二次写像が一般体で零なので全 $s$ で零であり、
全 $b\in k^4,c\in k^3$ が自由に取れる。定理\ref{thm:invariant-line-filtered-lifts}の
再構成により全てを実際の四テンソルに戻せ、閉点への良い還元で正則性が従う。

分解トーラスは階数1で、定理\ref{thm:split-spectral-line-classification}の軸交換が
指定分解を逆にする自己同型を与える。正負のせん断安定群は零なので恒等安定群はトーラスだけである。
従って無印の全流入が正方向で得られ、その非零点間の比較は向きを保つ群でよい。
この群の有限部分を、法線数 $n$ のテンソル上で以下のように作用させる。
$A,B$ は語中の文字0または2の個数についてそれぞれ $(-1)^{n_0},(-1)^{n_2}$ 倍、
$F$ は四位置の全ビット反転である。通常は $F^i A^j B^k$ の八個で尽きる。
実際、交代成分の $02|13$ 支持の二可分解線により全て単項で、
$L=1$ を保つ対角比較は二つの局所ビット比の二乗が1のものだけである。$s^4\ne1$ では
行列 $\left(\begin{smallmatrix}1&s\\s&1\end{smallmatrix}\right)$ の交差比から
独立な入替えは全反転だけである。$s^2=-1$ なら $s=\iota$ とし、さらに一つの作用 $E$ を加え、
八個の元と $E$ の左積を用いる。$E$ は各位置の文字2,3のビットだけを反転し、
旧語の各位置・文字に次表の倍率を掛ける。
\[
 \begin{pmatrix}
 -\iota&1&-1&-\iota\\
 \iota&1&\iota&1\\
 1&\iota&-\iota&1\\
 \iota&1&\iota&1
 \end{pmatrix}.
\]
これは28語の場合と同じ位置別対角比較を24支持語に適用するか、直接代入で確かめられる。
追加できる二つの独立な反転は $E,F$ で張られる。法線倍率は既に重み付き射影化に含まれる。
作用は三法線の接核をせん断列で割った $b$ 空間および四法線の接核 $c$ 空間にそのまま降ろす。
二次補正の切断をこのトーラスの有限拡大で同変に選ぶことで、この線型作用が全同型を与える。
\end{proof}
\begin{proof}[巡回の作用と全固有空間比較]
$S$ は左一巡回の後に文字2,3だけのビットを反転する。
旧語の位置0が文字2,3なら $s$、旧位置2が文字0,1なら $1/s$ を掛け、
最後に選んだ $r^2=s$ で $r^{n-2}$ 倍する。これは $n=2$ の背景を保ち、$S^4=1$ となる。
重み付き射影点 $[b,c]$ に周期 $m=1,2$ があるための必要十分条件は、前段の有限元 $g$ の一つで
\[
 b\ne0:\quad gS^m b=\lambda b,\quad gS^m c=\lambda^2 c;\qquad
 b=0:\quad c\text{ は }gS^m\text{ の固有ベクトル}
\]
となることである。二移動は全点で成立する。

計算を特殊値から全域へ戻す根拠を記す。$s^4\ne1$ では $r$ を変数とした連結な
開集合上で上の四元 $A,B,F,S$ が生成する作用は一つの有限2群を通る。
実際 $S$ の倍率は色だけによるので、$F$ と可換、$A$ と可換で $B$ を
法線倍率 $(-1)^{n-2}$ 倍に移す。$F$ も $A,B$ を共役で高々この中心符号だけ変える。
各生成元の位数は4以下である。$s=\iota$ では位相は八乗根、各位置内のビット反転と
巡回で作る置換も2群なので、ここでも全生成群が有限2群となる。
固有空間成分表は、二つの表現上のこの群の指標で決定される。
特に二つの固有空間の包含を調べるとき、その共通固有空間は対応する二作用の生成部分群の
一次元指標空間である。指定した二固有値が群関係に反すれば零であり、他は
指標冪等元で次元が決まる。これらは群の位数を逆にし必要な一の根を加えた
定階数格子上で直和因子になる。従って連結なパラメータ域で次元は一定であり、
良い奇素数への還元でも同じ包含表を得ることができる。

上記検算は通常行で $s=4,r=2$、例外行で $s=\iota$ を取り、法257ではそれぞれ
$(s,r)=(4,2),(16,4)$ とする。両接核とせん断商で次元を直接検査し、
各 $gS^m$ の全固有空間を取る。各語置換の位数は8以下、例外の位相は八乗根なので
六十四乗根までで足り、通常の生成関係でも十分である。法257にはこれらの根が全て存在する。
$b$ の固有空間と $c$ の固有値二乗空間の対、および $c$ のみの対から、
群作用の後に他の対に含まれるものを除く。$m=2$ では $k^4,k^3$ の対だけが残る。
$m=1$ では通常行の極大対の次元は $(2,3)$ 二つ、$(2,2)$ 四つ、$(2,1)$ 二つ、
例外行は $(2,3),(3,2),(2,2),(2,2),(2,1)$ である。
各々で射影化の一を引けば主張した成分を得る。
\end{proof}

\begin{corollary}\label{cor:all-split-attraction-atlas}
半安定商の閉軌道が可逆二平面分解を持つ全ての窓について、正則性、添字保持同型および
最小周期を決める有限な表と検査が得られる。この集合は正則特殊化で閉じ、
片側分岐中心を持つ正則窓とは交わらない。
\end{corollary}
\begin{proof}
閉軌道自身の候補は系\ref{cor:all-invertible-split-atlas}で尽くされる。
二固有値・単一固有値の非半単純直線は既出の量子または例外片側反転の閉軌道へ流入し、
$s=1$ の非半単純な射影合同の場合も二つの通常の行列のJordan成分を共役で消せば量子点になる。
重複曲面の非半単純型は反転あり半単純型へ流入する。従って新しい閉軌道を生じない。
量子基点には定理\ref{thm:quantum-universal-chambers}の全三表と安定群比較を用い、
一移動・二移動したテンソルに同じ有限比較を適用する。片側反転は
定理\ref{thm:mixed-flip-slice}, \ref{thm:mixed-flip-fiber}、
反転なし直線は上の定理、単純・重複曲面は定理\ref{thm:split-surface-all-fibers},
\ref{thm:reflected-collision-fibers}、六次式の孤立型は定理\ref{thm:isolated-split-component}である。
閉極限の一意性で異なる基点間の比較も終わる。特殊化閉性は定理\ref{thm:split-regular-limits}。
点対応が直線ファイバーを持つ条件は、点と隣の直線からなる射影パラメータ上で
全関係の評価消滅を課して射影した閉条件である。正則な閉極限では関係平面も階数6で特殊化する。
分解窓には定理\ref{thm:split-point-pencils}で左右ともそのようなファイバーがないため、
流入窓にも分岐はない。
\end{proof}
